% chapters/ch14-normalization-patterns.tex
\chapter{Normalization Patterns in Real Domains}

\begin{learningobjectives}
In this chapter you will learn a set of recurring normalization patterns used in real relational designs.
You will be able to recognise the underlying dependency behind each pattern, choose an appropriate schema structure, and explain what anomaly the pattern prevents.
\end{learningobjectives}

\section{Why patterns help}
After you learn the workflow and the core mapping rules, most design work becomes recognition.
You see a requirement and you match it to a known structure.
Patterns are not shortcuts for thinking.
They are a way to avoid repeating the same mistakes.

Each pattern below is presented with:
\begin{itemize}
  \item the situation in plain language,
  \item the dependency intuition,
  \item a compact schema in book notation,
  \item and the anomaly it prevents.
\end{itemize}

\section{Pattern 1: Order and OrderLine (the line-item pattern)}
\subsection{Situation}
A customer places an order.
An order contains multiple items.
Each item has a quantity and a price at the time of purchase.

The modelling trap is to store “items” as repeating groups inside \rel{Order}, or to store order attributes inside the line table.

\subsection{Dependency intuition}
Order facts are determined by \texttt{order\_id}.
Line facts are determined by (order\_id, line\_no) or by a surrogate \texttt{order\_line\_id}.
Product facts are determined by \texttt{product\_id}.
Do not mix them.

\subsection{Schema}
\[
\rel{Customer}(\underline{customer\_id}, name, email \;[\text{unique}])
\]
\[
\rel{Order}(\underline{order\_id}, created\_at, status, customer\_id \rightarrow \rel{Customer})
\]
\[
\rel{Product}(\underline{product\_id}, sku \;[\text{unique}], name, list\_price)
\]
\[
\rel{OrderLine}(\underline{order\_id} \rightarrow \rel{Order}, \underline{line\_no}, product\_id \rightarrow \rel{Product}, quantity, unit\_price)
\]

\subsection{What it prevents}
This structure prevents update anomalies like “customer name duplicated across order lines”.
It also supports history correctly: \texttt{unit\_price} is stored at purchase time, independent of future \texttt{list\_price} changes.

\section{Pattern 2: Address modelling (entity vs attribute)}
\subsection{Situation}
A system stores people or organisations and their addresses.
People can have multiple addresses (billing, shipping, current, previous).
Addresses can change over time.

A common failure is to place address fields directly on \rel{Customer}:
\texttt{street1, street2, city, zip, ...}.
That fails immediately when there is more than one address.

\subsection{Dependency intuition}
A customer does not determine a single address.
Instead, a customer determines a set of address relationships, possibly with types and validity periods.

\subsection{Schema (address as entity)}
\[
\rel{Address}(\underline{address\_id}, street, city, postal\_code, country)
\]
\[
\begin{aligned}
\rel{CustomerAddress}(&\underline{customer\_id} \rightarrow \rel{Customer}, \\
&\underline{address\_id} \rightarrow \rel{Address}, type, valid\_from, valid\_to?)
\end{aligned}
\]

\subsection{What it prevents}
This prevents repeating groups and allows history.
It also avoids mixing “what the address is” with “how the customer uses it” (billing/shipping), which belongs to the association.

\section{Pattern 3: Code and description (the lookup pattern)}
\subsection{Situation}
You store a code (status, category, specialty, role) and also store its human-readable name in multiple places.

This often creates inconsistent spellings and conflicting definitions.

\subsection{Dependency intuition}
If \texttt{status\_code} determines \texttt{status\_name}, then storing both together repeatedly creates redundancy.
The name belongs in one place.

\subsection{Schema}
\[
\rel{Status}(\underline{status\_code}, name, description)
\]
\[
\rel{Appointment}(\underline{appointment\_id}, start\_time, status\_code \rightarrow \rel{Status}, \dots)
\]

\subsection{What it prevents}
This prevents update anomalies where “cancelled” appears as “canceled” in some rows.
It also makes governance explicit: status definitions live in one relation.

\section{Pattern 4: Enrollment (association with attributes)}
\subsection{Situation}
A student can enroll in many courses.
A course has many students.
Enrollment has attributes like enrollment date, grade, and status.

\subsection{Dependency intuition}
Neither \texttt{student\_id} nor \texttt{course\_id} determines the other.
The association is the identity.
Enrollment attributes are determined by the pair.

\subsection{Schema}
\[
\rel{Student}(\underline{student\_id}, name, email)
\]
\[
\rel{Course}(\underline{course\_id}, code \;[\text{unique}], title)
\]
\[
\rel{Enrollment}(\underline{student\_id} \rightarrow \rel{Student}, \underline{course\_id} \rightarrow \rel{Course}, enrolled\_at, status, grade?)
\]

\subsection{What it prevents}
This prevents forced duplication such as storing a list of student ids inside a course row.
It also keeps enrollment history and supports constraints like “a student can be enrolled at most once in a course per term” (if you add \texttt{term}).

\section{Pattern 5: Inventory by location (context-dependent identity)}
\subsection{Situation}
A product exists across multiple locations (stores, warehouses).
You want quantity-on-hand per location.

A common mistake is to store one \texttt{quantity} on \rel{Product}.
That loses the location dimension.

\subsection{Dependency intuition}
Quantity is not determined by product alone.
It is determined by (product, location).

\subsection{Schema}
\[
\rel{Location}(\underline{location\_id}, name)
\]
\[
\rel{StockLevel}(\underline{product\_id} \rightarrow \rel{Product}, \underline{location\_id} \rightarrow \rel{Location}, quantity\_on\_hand)
\]

\subsection{What it prevents}
This prevents inconsistent inventory reporting and forces the model to represent location as a first-class dimension.

\section{Pattern 6: Audit log (append-only event history)}
\subsection{Situation}
You need a history of changes: who changed what, when, and from which value to which value.

A naive approach is to overwrite values and keep no record.
Another naive approach is to add many “previous\_*” columns.

\subsection{Dependency intuition}
History is not an attribute of the entity.
It is a stream of events, each with its own time and author.

\subsection{Schema}
\[
\rel{AuditEvent}(\underline{event\_id}, entity\_type, entity\_id, action, actor, occurred\_at)
\]
\[
\rel{AuditFieldChange}(\underline{event\_id} \rightarrow \rel{AuditEvent}, \underline{field\_name}, old\_value, new\_value)
\]

\subsection{What it prevents}
This avoids lossy overwrites and keeps traceability.
It also keeps audit structure separate from domain entities, preventing clutter and anomalies in core tables.

\section{How to apply patterns without forcing them}
Patterns are helpful, but they can be applied blindly.
The correct approach is:
\begin{itemize}
  \item start from domain meaning and rules,
  \item write a few dependencies in plain language,
  \item then choose the pattern that makes those dependencies align with keys.
\end{itemize}

When you choose a pattern, you should be able to say:
\emph{“This prevents this specific anomaly.”}
If you cannot, you may be restructuring for style rather than correctness.

\begin{chaptersummary}
Most relational schemas are combinations of a small set of patterns.
Line items separate order facts from item facts.
Addresses become entities because customers can have many addresses and address usage is a relationship.
Lookup tables centralize code meanings.
Association relations represent many-to-many membership and carry relationship attributes.
Context-dependent facts (like stock per location) belong to the pair.
Audit logs are append-only event history, separate from core entities.
\end{chaptersummary}

\begin{exercisebox}
Pick one domain you know and answer in full sentences, then write the schema fragments.
\begin{enumerate}[label=\arabic*.]
  \item Identify one place where you need a line-item pattern and write \rel{Order}/\rel{OrderLine}-style relations.
  \item Identify one attribute that is really a lookup code and design the lookup relation.
  \item Identify one many-to-many relationship that needs an association relation with attributes.
  \item Identify one context-dependent fact (product-by-location, member-by-branch, etc.) and model it.
  \item Choose one entity and propose a minimal audit logging structure for changes to it.
\end{enumerate}
\end{exercisebox}
